\documentclass[11pt,a4paper]{article}
\usepackage[utf8]{inputenc}
\usepackage[T1]{fontenc}
\usepackage{lmodern,amsmath,amssymb,amsthm,mathtools,booktabs,hyperref,enumitem,microtype}
\usepackage[margin=1in]{geometry}
\usepackage{natbib}
\newtheorem{definition}{Definition}[section]
\newtheorem{proposition}[definition]{Proposition}
\theoremstyle{remark}
\newtheorem{remark}[definition]{Remark}

\title{Adaptability\\\large One Operation Behind Shape Matching and Analogy, and Why It Is Not Called Intelligence}
\author{Matthew Parslow\\Independent Researcher}
\date{\textbf{Draft \textemdash\ draft-2026-08-30.4}\\[2pt]1 September 2026 \textbullet\ document revision 1}

\usepackage{fancyhdr}
\pagestyle{fancy}\fancyhf{}\fancyfoot[C]{\thepage}
\fancyfoot[R]{\scriptsize\textsc{draft-2026-08-30.4\,r1}}
\renewcommand{\headrulewidth}{0pt}
\fancypagestyle{plain}{\fancyhf{}\fancyfoot[C]{\thepage}%
  \fancyfoot[R]{\scriptsize\textsc{draft-2026-08-30.4\,r1}}\renewcommand{\headrulewidth}{0pt}}
\begin{document}
\maketitle
\begin{center}
\fbox{\parbox{0.92\textwidth}{\small
\textbf{DRAFT --- draft-2026-08-30.4 r1} (1 September 2026).
Circulated for comment, not as a finished result. This is the theory half of a pair; the companion
reports what is built and what is not. Sections marked as falsifiability commitments are the intended
points of attack.}}
\end{center}
\vspace{1em}

\begin{abstract}
Two capabilities this programme treats separately --- matching a shape against a held codebook, and
reasoning by analogy --- are one operation at two grains. Both \emph{extract the transform between
things already held and apply it where nothing is held}. Naming that operation gives a quantity worth
having: \emph{adaptability}, the capacity to acquire a competence that was not supplied. The word is
chosen against \emph{intelligence} deliberately. Intelligence is a scalar people rank; adaptability is
an operation with a cost, an ordering, and a failure mode, and a system either performs it on a given
case or does not.
\end{abstract}

\section{Why not intelligence}\label{sec:not-intelligence}
\emph{Intelligence} names a ranking before it names a mechanism. Asking whether a system has more of it
than another invites a scalar, and a scalar invites a benchmark, and a benchmark measures what a system
already does rather than what it can come to do. That is the wrong quantity for the question here.

\emph{Adaptability} is narrower and is an operation rather than a rank. A system is adaptable on a case
to the degree that, meeting something it does not hold, it ends up holding it --- and the interesting
part is the mechanism by which the new holding is \emph{derived from} the old rather than supplied. A
system given a competence has not adapted. A system that recognises everything it meets has not adapted
either; it was already sufficient. \emph{Adaptation requires a residual}, and the size of what must be
added to absorb one is the cost.

\section{The operation}\label{sec:operation}
\begin{definition}[Adaptation]\label{def:adapt}
Given a held set and an observation not covered by it, an \emph{adaptation} is the addition of the
smallest structure under which the observation is covered, derived from what is held.
\end{definition}

Three words carry it. \emph{Smallest}, because an adaptation that adds the observation itself is
storage, not adaptation. \emph{Derived}, because an adaptation supplied from outside is instruction.
And \emph{covered}, because the test is whether the case is now recognised, not whether a description
of it exists.

\begin{remark}[Shape matching and analogy are one operation]\label{rem:one}
Matching a shape against a codebook and completing an analogy look unrelated because their vocabularies
are. They are not.

To decode an observation against held signatures is to ask what transform carries a held one onto it.
Where that transform is the identity, the observation is recognised. Where it is a uniform step along a
dimension the codebook already has, the observation is a known thing moved --- and no new structure is
required. Where no such step exists but the residual is \emph{coupled} across the observed dimensions,
the transform requires a dimension not currently held, and one must be inferred before the observation
can be covered.

To complete an analogy is the same question with the arguments rearranged. Given two things held, the
transform between them is extracted; given a third, that transform is applied to reach a fourth which is
not held. \emph{Anti-unification --- generalising two structures while retaining what distinguishes
them --- is the general form, and the shape ladder is that form restricted}: at one rung the difference
is a single step along a known dimension, at the next it is a dimension.

So the two are not analogous to each other. They are the same operation, and the second is what the
first looks like when the objects are rules rather than observations.
\end{remark}

\section{The ladder, and why its order is not arbitrary}\label{sec:ladder}
An adaptation must be the smallest available, and that gives an ordering rather than a preference:

\begin{enumerate}[leftmargin=*]
\item \textbf{Recognise.} The observation is covered as held. Nothing is added, and no adaptation has
occurred --- which is the correct outcome for most observations.
\item \textbf{Transform in the held basis.} One uniform step along a dimension already carried. The
observation is a known thing displaced; the codebook does not grow.
\item \textbf{Infer a dimension.} The residual is coupled across observed dimensions in a way no step
in the current basis produces, so a dimension is minted and the transform is expressed in it.
\item \textbf{Mint.} No transform reaches the observation, so it is added as new.
\end{enumerate}

\emph{Only rungs three and four are adaptations}, and only the third is interesting: it is where the
system's own vocabulary changes. Rung four is the honest floor --- something genuinely new --- and a
system that reaches it too readily is not adapting but accumulating.

\begin{remark}[Why the order must be Occam and not preference]\label{rem:occam}
Attempting mint first would always succeed, so a system that did would never transform, never infer,
and would grow linearly in what it meets. Attempting transform first is what makes minting
\emph{informative}: a mint that survives the earlier rungs is evidence that nothing held reaches the
case, which is a claim about the held set and not only about the observation.

That is also why the ordering is a commitment and not an implementation detail. \emph{A system that
mints where a transform would have sufficed has recorded a coincidence as a category}, and one that
transforms where a mint was required has folded a distinction it needed. Both are failures of
adaptation, in opposite directions, and the ladder is what makes them separable.
\end{remark}

\section{What adaptability predicts}\label{sec:predicts}
The account is worth having only if it says what an adaptable system does that an unadaptable one does
not, on a case neither has seen.

\emph{One mint, then recognition.} Meeting a genuinely novel case repeatedly, an adaptable system should
mint \emph{exactly once} and recognise every later instance --- not mint per encounter, which is
storage, and not fail to mint, which is refusal. The signature is a step followed by silence.

\emph{Zero mints under transformation.} Meeting a known thing transformed along a held dimension, it
should add nothing at all. A system that mints there is treating the dimension it already has as absent.

\emph{Coupling forces dimensions.} A residual that is coupled across observed dimensions should produce
a dimension, not a member. \emph{The distinguishing case is the one whose difference cannot be written
in the current basis at all}, and a system without rung three cannot tell it apart from novelty --- it
mints where it should have generalised, and the codebook grows with the world instead of with its
structure.

\emph{And competence should transfer without being carried.} If the operation is real, structure built
for one domain should be usable unmodified in another whose existence was not anticipated when the
structure was made. That is the strongest available test, because it removes the author from the loop,
and it is the one the assessment companion reports against.

\section{What this does not claim}\label{sec:not-claimed}
Adaptability is not a quantity a system has more or less of in general. It is defined \emph{per case}
and against \emph{a held set}, so the same system is adaptable on one observation and not on another,
and moving the held set moves the verdict. There is no frame-free adaptability, for the same reason
this programme's base claim gives for there being no frame-free anything.

Nor does adaptation imply improvement. An adaptation is an addition that covers a residual; whether
covering it was worth the cost is a separate question with a separate answer, and a system that adapts
readily to noise is worse than one that refuses.

\section{Falsifiability commitments}\label{sec:falsifiability}
\begin{enumerate}[leftmargin=*]
\item \textsc{[shape]}~\textbf{Shape matching and analogy are one operation.} Fails if a case of one
cannot be expressed as a case of the other --- specifically, if some analogical completion has no
formulation as a transform between held signatures, or some shape adaptation has no formulation as an
anti-unification with its witness retained.
\item \textsc{[shape]}~\textbf{The ladder's order is forced, not chosen.} Fails if reordering the rungs
leaves the outcomes unchanged on cases where the rungs disagree, which would make the order a
preference.
\item \textsc{[detail]}~\textbf{One mint, then recognition.} Fails on a novel case that yields repeated
minting under an unchanged codebook, or none at all.
\item \textsc{[detail]}~\textbf{Coupling forces a dimension.} Fails if a residual coupled across
observed dimensions is absorbed by a member rather than a dimension, without loss.
\item \textsc{[shape]}~\textbf{Adaptability is per case, not a rank.} Fails if a single ordering of
systems by adaptability can be constructed that survives changing the held set.
\end{enumerate}

\end{document}
